\clearpage\chapter{Character theory}
\section*{1. Overview}
\subsection*{Intuitive}
\paragraph{The problem}

A representation turns each symmetry of a group into a matrix. Matrices can be large and change when the basis changes. Character theory compresses the representation to one function: the trace of each matrix. For finite groups over the complex numbers, this compact function contains enough information to recover the representation up to isomorphism.



\subsection*{Formal}
Let $\rho:G\to\operatorname{GL}(V)$ be a finite-dimensional representation of a finite group over a field of characteristic zero. Its character is the class function $\chi_\rho(g)=\operatorname{Tr}(\rho(g))$. The character inner product is
\[
\langle\chi,\psi\rangle=\frac1{|G|}\sum_{g\in G}\chi(g)\overline{\psi(g)}.
\]
Irreducible characters are orthonormal, and the multiplicity of an irreducible character in $\chi_\rho$ is its inner product with $\chi_\rho$.
\section*{2. History}
\subsection*{Historical development}
\begin{historylist}
\paragraph{Ferdinand Georg Frobenius}
Frobenius introduced characters as traces of representation matrices and used them to study finite-group representations without writing every matrix explicitly. In characteristic zero, the character determines a complex representation up to isomorphism, which makes character tables effective summary data.

\paragraph{Richard Brauer}
Brauer developed modular character theory for positive characteristic. His methods handle the case in which the characteristic divides the group order, where the averaging arguments used in ordinary character theory no longer apply.

These developments turned complicated group actions into numerical data that can be added, multiplied, and averaged. The resulting character methods support the study of finite groups, molecular symmetry, particle physics, and Lie groups \cite{character_theory_ref1,character_theory_ref3}.

\end{historylist}
\section*{3. Theory}
\subsection*{Core theory}
\paragraph{Representations and characters}

Let $\rho:G\to\operatorname{GL}(V)$ be a finite-dimensional representation over a field $F$. Its character is
\[
 \chi_\rho(g)=\operatorname{Tr}(\rho(g)).
\]
The degree is $\chi_\rho(1)=\dim V$ in characteristic zero. A character is irreducible when the representation has no nonzero proper invariant subspace. A character of degree one is linear.

Characters are class functions: $\chi(hgh^{-1})=\chi(g)$. The trace is unchanged by conjugating a matrix, so a character is constant on conjugacy classes. The kernel can be recovered in characteristic zero as the elements where $\chi(g)=\chi(1)$, and it is a normal subgroup. A character is not generally a group homomorphism.

\paragraph{Character tables}

For a finite group, make one column for each conjugacy class and one row for each irreducible character. The first row is the trivial character, whose entries are all 1; the first column is the identity class and contains the degrees.

The table is square because the number of irreducible complex representations equals the number of conjugacy classes. From the first column one obtains
\[
 |G|=\sum_i \chi_i(1)^2.
\]
The sum of squared absolute values in a column equals the order of the centralizer of a representative element. Kernels of characters reveal normal subgroups; the commutator subgroup is the intersection of kernels of linear characters. A character table can show whether a group is abelian, but it need not determine the group up to isomorphism: the quaternion group and the dihedral group of order eight have the same character table.

\paragraph{Orthogonality and decomposition}

For complex class functions define
\[
 \langle\alpha,\beta\rangle=
 \frac1{|G|}\sum_{g\in G}\alpha(g)\overline{\beta(g)}.
\]
Irreducible characters form an orthonormal basis, so
\[
 \langle\chi_i,\chi_j\rangle=\delta_{ij}.
\]
If $\chi=\sum_i m_i\chi_i$, then $m_i=\langle\chi,\chi_i\rangle$. Thus the inner product tells how many times each irreducible representation occurs in a decomposition.

The column orthogonality relation says that the corresponding columns are orthogonal after weighting by centralizers. These relations are practical algorithms for completing a character table, finding centralizer sizes, decomposing a representation, and checking whether a proposed row is possible.

\paragraph{Arithmetic rules}

Characters turn common representation operations into simple formulas:
\[
 \chi_{\rho\oplus\sigma}=\chi_\rho+\chi_\sigma,\qquad
 \chi_{\rho\otimes\sigma}=\chi_\rho\chi_\sigma,\qquad
 \chi_{\rho^*}=\overline{\chi_\rho}.
\]
For exterior and symmetric squares,
\[
 \chi_{\wedge^2\rho}(g)=\frac{\chi_\rho(g)^2-\chi_\rho(g^2)}2,
\quad
 \chi_{\operatorname{Sym}^2\rho}(g)=\frac{\chi_\rho(g)^2+\chi_\rho(g^2)}2.
\]
For finite-order $g$, the eigenvalues of $\rho(g)$ are roots of unity, so character values are sums of roots of unity and hence algebraic integers over $\mathbb C$.

\paragraph{Restriction, induction, and Frobenius reciprocity}

If $H\leq G$, restricting a representation from $G$ to $H$ restricts its character. Conversely, a representation of $H$ can be induced to a representation of $G$ by letting $G$ act on functions or cosets with values in the original space. The induced character is zero on elements not conjugate to an element of $H$ and has a coset-sum formula on elements of $H$.

Frobenius reciprocity says
\[
 \langle\theta^G,\chi\rangle_G
 =\langle\theta,\chi_H\rangle_H.
\]
It converts multiplicity calculations for induction into calculations for restriction and is one of the main tools for building character tables from subgroups.

\paragraph{Lie groups and Lie algebras}

For a finite-dimensional representation of a Lie group, the character is still $\operatorname{Tr}(\rho(g))$. For a Lie algebra representation, one may define
\[
 \chi_\rho(X)=\operatorname{Tr}(e^{\rho(X)}).
\]
The character is invariant under the adjoint action. For a complex semisimple Lie algebra, it is determined by its values on a Cartan subalgebra and can be written using weights:
\[
 \chi_\rho(H)=\sum_\lambda m_\lambda e^{\lambda(H)}.
\]
The Weyl character formula computes irreducible characters explicitly from roots and weights.

\section*{4. Important theorems and results}
\begin{enumerate}[leftmargin=*,itemsep=0.35em]

\item \textbf{Character determination theorem.} Complex representations of a finite group are isomorphic exactly when they have the same character.

\item \textbf{Schur orthogonality.} Irreducible characters form an orthonormal basis of the class functions, giving both row and column orthogonality.

\item \textbf{Sum-of-squares theorem.} The order of a finite group equals the sum of squares of the degrees of its irreducible complex representations.

\item \textbf{Maschke's theorem.} If the characteristic does not divide $|G|$, every finite-dimensional representation is completely reducible. This is why character decomposition works cleanly in characteristic zero.

\item \textbf{Frobenius reciprocity.} Multiplicities in an induced representation equal multiplicities in the corresponding restricted representation.

\item \textbf{Weyl character formula.} The character of an irreducible representation of a complex semisimple Lie algebra is determined explicitly by its highest weight and root system.
\end{enumerate}
\noindent\emph{Source for these results:} \cite{character_theory_ref1}.

\section*{5. Applications}
Character theory compresses a representation into traces of group elements, while retaining enough information to study its decomposition. Character values therefore make symmetry calculations manageable in algebra, combinatorics, chemistry, and physics.
\begin{enumerate}[leftmargin=*,itemsep=0.35em]
\item \textbf{Finite-group structure.} Character values reveal normal subgroups, commutators, centralizers, and possible degrees. Character calculations occur in classification of finite simple groups and in the Feit--Thompson theorem.
\item \textbf{Chemistry and physics.} Molecular vibrations and particle states transform under symmetry groups. Character tables separate modes into irreducible types and predict which transitions or interactions are allowed.
\item \textbf{Combinatorics.} The representation of a group on a set gives a permutation character. Inner products count orbits and fixed configurations, connecting characters with Burnside's lemma.
\item \textbf{Number theory and analysis.} Linear characters of abelian groups include Dirichlet characters and lead to Fourier analysis and $L$-functions. Higher-dimensional characters enter automorphic forms and arithmetic geometry.
\item \textbf{Moonshine.} Replacing ordinary dimension by a graded character produces McKay--Thompson series and connects representations of the Monster group with modular functions.
\end{enumerate}


\theoryconnections{character_theory}{%
\item Group theory provides conjugacy classes, which collect elements that have the same symmetry type.
\item Representation theory turns each group element into a matrix, and the character records only the trace of that matrix; traces are easier to compare and still remember enough to decompose a representation.
\item Orthogonality of characters is a Fourier calculation on the group, so inner products count how many irreducible symmetry patterns occur \cite{character_theory_ref1,character_theory_ref2}.
}{%
\item Character theory helps classify finite-group representations by replacing a large collection of matrices with a table of values on conjugacy classes.
\item In chemistry those values predict how molecular vibrations split into symmetry types; in particle physics they label states under continuous symmetries.
\item In combinatorics and number theory, character sums convert symmetry or congruence conditions into cancellations that can be estimated \cite{character_theory_ref1,character_theory_ref3}.
}

\section*{Glossary}
\begin{description}[leftmargin=*,style=nextline,itemsep=0.3em]

\item[\textbf{Character.}] The class function $\chi_\rho(g)=\operatorname{Tr}(\rho(g))$ assigning to each group element the trace of the corresponding representation matrix; it determines the representation up to isomorphism for finite groups over $\mathbb C$.

\item[\textbf{Irreducible character.}] The character of a representation that has no nonzero proper invariant subspaces; irreducible characters form an orthonormal basis for class functions.

\item[\textbf{Conjugacy class.}] A set of group elements that are all conjugate to one another; characters are constant on conjugacy classes because the trace is unchanged by similarity.

\item[\textbf{Character table.}] A square matrix whose rows are irreducible characters and whose columns are conjugacy classes, serving as a complete summary of a finite group's representation theory.

\item[\textbf{Induced character.}] The character of a representation of a larger group $G$ constructed from a representation of a subgroup $H$ by letting $G$ act on coset copies of the original space.

\item[\textbf{Frobenius reciprocity.}] A duality stating that the multiplicity of an irreducible character of $G$ in an induced representation equals the multiplicity of the original character of $H$ in the restriction.

\item[\textbf{Class function.}] A function on a group that is constant on each conjugacy class; the space of class functions is the natural domain for character inner products.

\item[\textbf{Representation.}] A group homomorphism $\rho:G\to\operatorname{GL}(V)$ realizing each group element as a linear transformation; characters are traces of these matrices.

\item[\textbf{Maschke's theorem.}] Every finite-dimensional representation of a finite group is completely reducible when the field characteristic does not divide the group order.

\item[\textbf{Schur orthogonality.}] The relations stating irreducible characters form an orthonormal basis for class functions.

\item[\textbf{Degree of a character.}] The dimension $\chi(1)$ of the underlying representation space; it appears in the sum-of-squares formula.

\end{description}

\section*{7. Sample Questions}
\emph{Exercise reference:} \cite{character_theory_ref1}. The cited edition is the source for the exercise set; consult its Exercises section for the official statement and numbering.
\begin{enumerate}[leftmargin=*,itemsep=0.9em]

\item \textbf{Exercise 1.} Construct the character table of $\mathbb{Z}/4\mathbb{Z}$.

\emph{Solution.} $\mathbb{Z}/4\mathbb{Z}$ is abelian with 4 conjugacy classes (each element is its own class) and 4 irreducible characters, all of degree 1. Let $\omega=i=\sqrt{-1}$. The characters are $\chi_k(g)=\omega^{kg}$ for $k=0,1,2,3$ and $g=0,1,2,3$:

\[
\begin{array}{c|cccc}
& 0 & 1 & 2 & 3\\\hline
\chi_0 & 1 & 1 & 1 & 1\\
\chi_1 & 1 & i & -1 & -i\\
\chi_2 & 1 & -1 & 1 & -1\\
\chi_3 & 1 & -i & -1 & i
\end{array}
\]

Check: $\sum_k|\chi_k(1)|^2=4=|\mathbb{Z}/4\mathbb{Z}|$.

\item \textbf{Exercise 2.} Let $\chi$ be the character of the standard 2-dimensional representation of $S_3$ (the representation on $\{(x,y,z):x+y+z=0\}$). Compute $\langle\chi,\chi\rangle$ and determine whether the representation is irreducible.

\emph{Solution.} The conjugacy classes of $S_3$ are: identity (size 1), transpositions (size 3), 3-cycles (size 2). The character values are $\chi(1)=2$, $\chi((12))=0$, $\chi((123))=-1$. Then $\langle\chi,\chi\rangle=\frac{1}{6}(1\cdot|2|^2+3\cdot|0|^2+2\cdot|-1|^2)=\frac{1}{6}(4+0+2)=1$. Since $\langle\chi,\chi\rangle=1$, the representation is irreducible.

\item \textbf{Exercise 3.} Compute the character of the tensor product $\chi\otimes\psi$ where $\chi$ and $\psi$ are the two nontrivial irreducible characters of $\mathbb{Z}/3\mathbb{Z}$, evaluated at the generator $g$ of order 3.

\emph{Solution.} Let $\omega=e^{2\pi i/3}$. The irreducible characters of $\mathbb{Z}/3\mathbb{Z}=\{0,1,2\}$ are $\chi_0(g^k)=1$, $\chi_1(g^k)=\omega^k$, $\chi_2(g^k)=\omega^{2k}$. The nontrivial ones are $\chi_1$ and $\chi_2$. Their tensor product: $(\chi_1\otimes\chi_2)(g^k)=\omega^k\cdot\omega^{2k}=\omega^{3k}=1$. So $\chi_1\otimes\chi_2=\chi_0$, the trivial character. This reflects the fact that $\chi_2=\overline{\chi_1}$ and $\chi_1\cdot\overline{\chi_1}=|\chi_1|^2=1$.

\item \textbf{Exercise 4.} The character of the regular representation of a finite group $G$ is $\chi_{\text{reg}}(g)=|G|$ if $g=e$ and $0$ otherwise. Decompose $\chi_{\text{reg}}$ into irreducibles for $G=\mathbb{Z}/2\mathbb{Z}$.

\emph{Solution.} $G=\{e,g\}$, $|G|=2$. The irreducible characters are $\chi_0=(1,1)$ (trivial) and $\chi_1=(1,-1)$ (sign). The regular character is $\chi_{\text{reg}}=(2,0)$. Then $m_0=\langle\chi_{\text{reg}},\chi_0\rangle=\frac{1}{2}(2\cdot 1+0\cdot 1)=1$ and $m_1=\langle\chi_{\text{reg}},\chi_1\rangle=\frac{1}{2}(2\cdot 1+0\cdot(-1))=1$. So $\chi_{\text{reg}}=\chi_0+\chi_1$, and the regular representation decomposes as the direct sum of both irreducibles, each appearing once.

\item \textbf{Exercise 5.} Use the sum-of-squares formula to determine the order of a group whose irreducible character degrees are $1,1,1,2$.

\emph{Solution.} By the formula $|G|=\sum_i\chi_i(1)^2=1^2+1^2+1^2+2^2=1+1+1+4=7$. So $|G|=7$. Since 7 is prime, $G\cong\mathbb{Z}/7\mathbb{Z}$, and indeed this abelian group has three degree-1 characters... wait, $\mathbb{Z}/7\mathbb{Z}$ has 7 conjugacy classes and 7 irreducible characters, all of degree 1. The degrees $1,1,1,2$ sum to $|G|=7$ but require 4 conjugacy classes. In fact $|G|=1+1+1+4=7$ works for a nonabelian group of order 7? No—order 7 is abelian. Let me restate: there is no group with exactly these degrees since $|G|=7$ forces all degrees to be 1. The correct interpretation: the sum-of-squares formula gives $|G|=7$, but this is inconsistent with having a degree-2 representation (which requires $|G|\ge 6$ and nonabelian, but $|G|=7$ is abelian). So no such group exists. The formula $|G|=7$ is the answer, and the inconsistency shows these cannot be the character degrees of any group.

\end{enumerate}
\section*{8. References}
\begin{thebibliography}{99}
\bibitem{character_theory_ref1} J.-P. Serre, \emph{Linear Representations of Finite Groups}, Springer, 1977.
\bibitem{character_theory_ref2} I. M. Isaacs, \emph{Character Theory of Finite Groups}, Academic Press, 1976.
\bibitem{character_theory_ref3} B. C. Hall, \emph{Lie Groups, Lie Algebras, and Representations}, 2nd edition, Springer, 2015.
\end{thebibliography}
